\subsubsection{闭算子}\label{subsec:closed-operators}

以下概念在研究无界算子时非常有用.

\begin{Definition}
\label{def:linear-operator-graph}
线性算子 $(\mathcal{D}(A),A)$ 的\textbf{图}是如下有序对的集合:
\begin{equation}\label{eq:linear-operator-graph-definition}
\Gamma(A):=\{(x,Ax)\mid x\in\mathcal{D}(A)\}\subset X\times Y.
\end{equation}

\end{Definition}

注意, 图是 $X\times Y$ 的子空间.

以下引理是算子的扩张和图的定义的直接推论.

\begin{Lemma}
\label{lem:operator-extension-graph-inclusion}
算子 $(\mathcal{D}(\widetilde A),\widetilde A)$ 是 $(\mathcal{D}(A),A)$ 的扩张, 当且仅当 $\Gamma(\widetilde A)\supset\Gamma(A)$.
\end{Lemma}

证明留给读者.

\begin{Definition}
\label{def:closed-closable-operator-and-closure}
若算子 $(\mathcal{D}(A),A)$ 的图作为 $X\times Y$ 的子集是闭的, 则称该算子是\textbf{闭的}. 若 $(\mathcal{D}(A),A)$ 有闭扩张, 则称它是\textbf{可闭的}. 每个可闭算子都有一个最小闭扩张, 称为它的\textbf{闭包}, 记为 $(\overline{\mathcal{D}(A)},\overline A)$.
\end{Definition}

用闭算子的 ``序列'' 概念来补充这一定义是有用的. 例如, 以下引理是闭算子以及乘积空间中闭集的定义的直接推论.

\begin{Lemma}
\label{lem:closed-operator-sequential-criterion}
算子 $(\mathcal{D}(A),A)$ 是闭的, 当且仅当它具有以下性质. 只要存在序列 $x_n\in\mathcal{D}(A)$ 满足\begin{enumerate}\item $x_n\to x$, 且\item $Ax_n\to f$,\end{enumerate}就有\begin{enumerate}\item $x\in\mathcal{D}(A)$, 且\item $Ax=f$.\end{enumerate}
\end{Lemma}

对于可闭算子, 有类似的刻画.

\begin{Lemma}
\label{lem:closable-operator-zero-sequence-criterion}
若对于每个满足 $x_n\to0$ 的序列 $x_n\in\mathcal{D}(A)$, 均有以下两者之一:\begin{enumerate}\item $Ax_n\to0$, 或者\item $\lim_{n\to\infty}Ax_n$ 不存在.\end{enumerate}则算子 $(\mathcal{D}(A),A)$ 可闭.
\end{Lemma}

\begingroup
\postdisplaypenalty=10000
\renewcommand{\endProof}{\par\nobreak\hfill$\blacksquare$\par\vspace{0.1cm}}
\begin{Proof}
为证明这一点, 如下构造 $(\mathcal{D}(A),A)$ 的闭扩张 $(\mathcal{D}(\widetilde A),\widetilde A)$. 若 $x_n\in\mathcal{D}(A)$ 是满足以下条件的序列:\begin{enumerate}\item $x_n\to x$, 且\item $Ax_n\to f$,\end{enumerate}其中 $f\in Y$, 则令\begin{enumerate}\item $x\in\mathcal{D}(\widetilde A)$, 并定义\item $\widetilde Ax=f$.\end{enumerate}需要确保 $\widetilde Ax$ 的定义没有歧义. 然而, 假设保证, 若另有序列 $\widehat x_n\in\mathcal{D}(A)$ 且 $\widehat x_n\to x$, 则有以下两者之一:\begin{enumerate}\item $A\widehat x_n\to f$, 或者\item $\lim_{n\to\infty}A\widehat x_n$ 不存在\end{enumerate}(因为 $x_n-\widehat x_n\to0$). 根据 Lemma~\ref{lem:closed-operator-sequential-criterion}, 算子 $(\mathcal{D}(\widetilde A),\widetilde A)$ 是闭的, 而且按其构造, 它是 $(\mathcal{D}(A),A)$ 的扩张.
\end{Proof}
\endgroup

事实上, 容易看出 $\Gamma(\widetilde A)=\overline{\Gamma(A)}$. 因此, 由于 $\overline{\Gamma(A)}\subset\Gamma(\overline A)$ (这直接来自 Definition~\ref{def:closed-closable-operator-and-closure}), 而且闭包按定义是尽可能小的扩张, 上述构造实际上给出了闭包; 即 $(\mathcal{D}(\widetilde A),\widetilde A)=(\overline{\mathcal{D}(A)},\overline A)$. 事实上, 已经证明以下结果.

\begin{Corollary}
\label{cor:closure-graph-is-graph-closure}
若 $(\mathcal{D}(A),A)$ 可闭, 则 $\Gamma(\overline A)=\overline{\Gamma(A)}$.
\end{Corollary}
